\documentclass[11pt,a4paper]{article}

% ----------------------------------------------------------------------------
%  Codificação e idioma
% ----------------------------------------------------------------------------
\usepackage[utf8]{inputenc}
\usepackage[T1]{fontenc}
\usepackage[brazil]{babel}
\usepackage{lmodern}

% ----------------------------------------------------------------------------
%  Layout
% ----------------------------------------------------------------------------
\usepackage[a4paper,margin=2.5cm]{geometry}
\usepackage{setspace}
\onehalfspacing
\usepackage{parskip}

% ----------------------------------------------------------------------------
%  Matemática, tabelas, cores
% ----------------------------------------------------------------------------
\usepackage{amsmath,amssymb}
\usepackage{booktabs}
\usepackage{array}
\usepackage[table]{xcolor}
\usepackage{caption}
\usepackage{subcaption}

% ----------------------------------------------------------------------------
%  Algoritmos e código
% ----------------------------------------------------------------------------
\usepackage{algorithm}
\usepackage{algpseudocode}
\usepackage{listings}

% ----------------------------------------------------------------------------
%  Diagramas
% ----------------------------------------------------------------------------
\usepackage{tikz}
\usetikzlibrary{arrows.meta,positioning,shapes.geometric,fit,backgrounds}

% ----------------------------------------------------------------------------
%  Links (por último)
% ----------------------------------------------------------------------------
\usepackage[colorlinks=true,linkcolor=blue!60!black,
            citecolor=green!50!black,urlcolor=blue!60!black]{hyperref}

% ----------------------------------------------------------------------------
%  Cores e estilo de código
% ----------------------------------------------------------------------------
\definecolor{cellchg}{RGB}{198,239,206}   % célula alterada (verde)
\definecolor{cellrcl}{RGB}{197,217,241}   % elemento na RCL (azul)
\definecolor{cellsel}{RGB}{252,213,180}   % elemento escolhido (laranja)
\definecolor{codebg}{RGB}{248,248,248}
\definecolor{codekw}{RGB}{0,0,180}
\definecolor{codecm}{RGB}{0,128,0}
\definecolor{codestr}{RGB}{170,55,55}

\lstdefinestyle{py}{
  language=Python,
  backgroundcolor=\color{codebg},
  basicstyle=\ttfamily\footnotesize,
  keywordstyle=\color{codekw}\bfseries,
  commentstyle=\color{codecm}\itshape,
  stringstyle=\color{codestr},
  numbers=left,
  numberstyle=\tiny\color{gray},
  numbersep=8pt,
  showstringspaces=false,
  breaklines=true,
  frame=single,
  rulecolor=\color{gray!40},
  tabsize=4,
  captionpos=b,
  keepspaces=true,
}

% Comando auxiliar: rótulo de hexágono
\newcommand{\hex}[1]{\textsf{#1}}

% ----------------------------------------------------------------------------
%  Capa
% ----------------------------------------------------------------------------
\title{\textbf{Implementação do GRASP para Otimização do Índice de\\
Qualidade de Caminhabilidade (IQC)}\\[4pt]
\large Relatório técnico detalhado: teoria e implementação passo a passo}
\author{Anderson Oliveira}
\date{\today}

\begin{document}
\maketitle

\begin{abstract}
\noindent
Este relatório documenta a implementação do \emph{Greedy Randomized Adaptive
Search Procedure} (GRASP), na sua forma clássica de duas fases descrita por
Resende e Ribeiro, aplicado ao problema de alocação de Pontos de Interesse
(POIs) que maximiza o Índice de Qualidade de Caminhabilidade (IQC) total de uma
malha de hexágonos H3. A primeira parte apresenta a teoria do GRASP utilizada
(multi-start, construção gulosa-aleatória com Lista Restrita de Candidatos e
busca local). A segunda parte descreve, função a função, como o método foi
implementado, incluindo as estruturas de dados em \texttt{ndarray}, a função
objetivo (pesos CRITIC + IQC) e as decisões de projeto. Por fim, uma instância
didática é executada e \emph{ilustrada passo a passo}, mostrando como a matriz de
alocação evolui, como o impacto se propaga entre hexágonos e como o IQC é
recalculado a cada iteração.
\end{abstract}

\tableofcontents
\newpage

% ============================================================================
\section{Introdução e contexto}
% ============================================================================

O projeto calcula, para cada hexágono H3 de uma localidade, um conjunto de
indicadores de caminhabilidade (saúde, educação, transporte, segurança, etc.) e
deles deriva o \textbf{IQC} (Índice de Qualidade de Caminhabilidade) via pesos
\textbf{CRITIC}. A etapa de otimização busca responder a uma pergunta de
planejamento urbano:

\begin{quote}
\itshape
Dado um orçamento de $B$ novos POIs, em quais hexágonos e de quais dimensões
devemos inseri-los para maximizar o IQC \emph{total} da região?
\end{quote}

Trata-se de um problema combinatório: o número de maneiras de distribuir $B$
POIs entre os pares (hexágono-fonte, dimensão) cresce explosivamente. Métodos
exatos são inviáveis para instâncias reais, o que motiva o uso de uma
\textbf{meta-heurística}. Escolhemos o GRASP por três razões alinhadas ao texto
de Resende e Ribeiro: (i) requer poucos parâmetros (essencialmente o parâmetro
$\alpha$ da lista de candidatos); (ii) tem estrutura simples de implementar sobre
blocos já disponíveis (uma função de construção e uma de busca local); e (iii)
é naturalmente multi-start, o que o torna robusto à aleatoriedade das sementes.

% ============================================================================
\section{Formalização do problema}
\label{sec:problema}
% ============================================================================

\subsection{Conjunto base e variável de decisão}

Seja $H=\{h_1,\dots,h_n\}$ o conjunto dos $n$ hexágonos e
$D=\{d_1,\dots,d_m\}$ o conjunto das $m$ dimensões de POI alocáveis (saúde,
transporte, etc.). O \textbf{conjunto base} (\emph{ground set}) do GRASP é

\begin{equation}
E \;=\; H_{\text{fonte}} \times D,
\end{equation}

\noindent
onde $H_{\text{fonte}}\subseteq H$ são os hexágonos elegíveis para receber POIs.
Cada \emph{elemento} $e=(h,d)\in E$ representa ``inserir uma unidade de POI da
dimensão $d$ no hexágono-fonte $h$''.

A solução é uma \textbf{matriz de alocação} (chamada no código de
\texttt{candidate\_matrix})
\begin{equation}
X \in \mathbb{Z}_{\ge 0}^{\,n \times m},
\qquad
X_{h,d} = \text{número de POIs da dimensão } d \text{ alocados no hexágono } h .
\end{equation}

\subsection{Restrição de viabilidade}

A única restrição é o orçamento total:
\begin{equation}
\sum_{h=1}^{n}\sum_{d=1}^{m} X_{h,d} \;=\; B .
\label{eq:budget}
\end{equation}
Como a construção do GRASP insere exatamente um POI por passo e para em $B$
passos, e a busca local apenas \emph{move} POIs (preservando a soma), a
restrição~\eqref{eq:budget} é satisfeita por construção --- nenhum procedimento
de reparo é necessário.

\subsection{Propagação espaço-temporal do impacto}

Inserir um POI num hexágono-fonte $s$ não beneficia apenas $s$: o efeito
se \emph{propaga} aos hexágonos-alvo $t$ alcançáveis a pé, ponderado por um fator
de decaimento temporal $\alpha_{s,t}\in(0,1]$ (cosseno do tempo de caminhada,
campo \texttt{alpha\_20}). Seja $P$ o conjunto de pares fonte--alvo com
$\alpha_{s,t}>0$. A matriz final de indicadores é
\begin{equation}
M \;=\; M^{0} + \Delta,
\qquad
\Delta_{t,\,j(d)} \;=\; \sum_{(s,t)\in P} \alpha_{s,t}\, X_{s,d},
\label{eq:propagacao}
\end{equation}
onde $M^{0}$ é a matriz \emph{baseline} de indicadores ($n\times 11$),
$j(d)$ é a coluna do indicador correspondente à dimensão $d$, e a soma
acumula a contribuição de todos os POIs que, a partir de suas fontes, alcançam
o alvo $t$.

\subsection{Função objetivo: CRITIC + IQC}

Dada a matriz final $M$, a função objetivo executa três etapas
(implementadas em \texttt{objective\_function}):

\paragraph{(1) Pesos CRITIC.} Normaliza-se cada coluna por \emph{min--max},
\begin{equation}
\hat{M}_{ij} = \frac{M_{ij}-\min_i M_{ij}}{\max_i M_{ij}-\min_i M_{ij}},
\end{equation}
calcula-se o desvio-padrão $\sigma_j$ de cada coluna normalizada e o
\emph{conflito} $c_j=\sum_k\bigl(1-|\rho_{jk}|\bigr)$, onde $\rho_{jk}$ é a
correlação de Pearson entre os indicadores $j$ e $k$. O peso de cada indicador é
\begin{equation}
w_j \;=\; \frac{\sigma_j\, c_j}{\sum_{l}\sigma_l\, c_l}.
\label{eq:critic}
\end{equation}

\paragraph{(2) IQC por hexágono.} É a média ponderada dos indicadores
normalizados:
\begin{equation}
\mathrm{IQC}_i \;=\; \sum_{j} w_j\, \hat{M}_{ij}.
\label{eq:iqc}
\end{equation}

\paragraph{(3) Objetivo global.} Soma do IQC sobre todos os hexágonos, a ser
\textbf{maximizada}:
\begin{equation}
\boxed{\;f(X) \;=\; \sum_{i=1}^{n}\mathrm{IQC}_i\;}
\qquad\text{(direção: maximizar).}
\label{eq:objetivo}
\end{equation}

\begin{quote}\small
\textbf{Observação crucial.} Os pesos CRITIC~\eqref{eq:critic} e a
normalização~\eqref{eq:iqc} dependem de \emph{todas} as colunas e linhas de $M$.
Logo, $f$ é \textbf{não-separável}: o ganho de inserir um POI não é constante ---
muda conforme a solução parcial atual. Isso encaixa exatamente no caráter
\emph{adaptativo} do GRASP (a função gulosa é reavaliada a cada passo), mas torna
cada avaliação cara, o que orienta as decisões de eficiência da
Seção~\ref{sec:impl}.
\end{quote}

% ============================================================================
\section{Teoria do GRASP utilizada}
\label{sec:teoria}
% ============================================================================

O GRASP (Feo e Resende, 1995; Resende e Ribeiro, 2016, 2019) é uma
meta-heurística \emph{multi-start}: cada iteração é independente e consiste em
duas fases --- \textbf{construção} gulosa-aleatória e \textbf{busca local} --- e a
melhor solução global é mantida ao final.

\begin{figure}[ht]
\centering
\begin{tikzpicture}[
  node distance=8mm and 14mm,
  box/.style={rectangle,rounded corners,draw=black!60,fill=blue!8,
              minimum width=3.2cm,minimum height=1cm,align=center},
  arr/.style={-{Stealth[length=2.5mm]},thick}]
  \node[box] (seed) {Para cada\\semente $\sigma$};
  \node[box,right=of seed,fill=green!10] (cons) {Construção\\gulosa-aleatória\\(RCL, $\alpha$)};
  \node[box,right=of cons,fill=orange!12] (ls) {Busca local\\(first-improving)};
  \node[box,right=of ls,fill=red!8] (upd) {Atualiza\\melhor solução $S^\star$};
  \draw[arr] (seed) -- (cons);
  \draw[arr] (cons) -- node[above,font=\scriptsize]{$S$} (ls);
  \draw[arr] (ls) -- node[above,font=\scriptsize]{$S$ ótimo local} (upd);
  \draw[arr] (upd.south) .. controls +(0,-1.1) and +(0,-1.1) .. (seed.south)
        node[midway,below,font=\scriptsize]{próxima semente (multi-start)};
\end{tikzpicture}
\caption{Estrutura geral do GRASP: laço multi-start sobre as sementes;
cada iteração constrói uma solução, refina-a por busca local e atualiza a melhor
solução encontrada.}
\label{fig:grasp-flow}
\end{figure}

\subsection{Visão geral (multi-start)}

O Algoritmo~\ref{alg:grasp} resume o laço principal, adaptado para
\textbf{maximização} (o texto original trata minimização). \texttt{Max\_Iterations}
corresponde, na nossa implementação, ao número de sementes disponíveis.

\begin{algorithm}[ht]
\caption{GRASP (versão de maximização)}
\label{alg:grasp}
\begin{algorithmic}[1]
\State $f^\star \gets -\infty$;\quad $S^\star \gets \varnothing$
\ForAll{semente $\sigma$ em \textit{seeds}}
  \State $S \gets \textsc{ConstrucaoGulosaAleatoria}(\alpha,\sigma)$
  \State $S \gets \textsc{BuscaLocal}(S)$
  \If{$f(S) > f^\star$}
    \State $S^\star \gets S$;\quad $f^\star \gets f(S)$
  \EndIf
\EndFor
\State \Return $S^\star$
\end{algorithmic}
\end{algorithm}

\subsection{Fase de construção e a Lista Restrita de Candidatos (RCL)}

A construção monta a solução \emph{um elemento por vez}. A cada passo
avalia-se, para cada elemento candidato $e$, o \textbf{ganho incremental}
$g(e)$ que sua inserção provoca na função objetivo. Constrói-se então a
\textbf{Lista Restrita de Candidatos} (RCL) com os \emph{melhores} elementos e
sorteia-se um deles ao acaso. Três aspectos coexistem:

\begin{itemize}
\item \textbf{Guloso} --- só entram na RCL os elementos de maior ganho;
\item \textbf{Probabilístico} --- escolhe-se aleatoriamente \emph{dentro} da RCL;
\item \textbf{Adaptativo} --- após cada inserção os ganhos são reavaliados.
\end{itemize}

A RCL é controlada pelo limiar $\alpha\in[0,1]$. Para \textbf{maximização}, sendo
$g_{\max}$ e $g_{\min}$ o maior e o menor ganho dos candidatos, a lista é
\begin{equation}
\mathrm{RCL} \;=\;\Bigl\{\, e : g(e) \;\ge\; g_{\max} - \alpha\,(g_{\max}-g_{\min}) \,\Bigr\}.
\label{eq:rcl}
\end{equation}
Os extremos têm interpretação direta: $\alpha=0$ produz um algoritmo
\emph{puramente guloso} (só o melhor elemento entra na RCL), enquanto $\alpha=1$
produz uma \emph{construção aleatória} (todos entram). Resende e Ribeiro mostram
empiricamente que valores intermediários ($\alpha\approx 0{,}2$) oferecem o melhor
compromisso entre média e variância das soluções; adotamos esse valor.

\begin{algorithm}[ht]
\caption{Construção gulosa-aleatória (maximização, \emph{sampled greedy})}
\label{alg:cons}
\begin{algorithmic}[1]
\State $S \gets$ matriz nula $n\times m$;\quad $v \gets f(S)$
\For{$k = 1,\dots,B$}
  \State $\mathcal{C} \gets$ amostra aleatória de $p$ pares $(h,d)$ do conjunto base $E$
  \ForAll{$(h,d) \in \mathcal{C}$}
    \State $g(h,d) \gets f\bigl(S \text{ com } X_{h,d}{+}1\bigr) - v$ \Comment{ganho incremental}
  \EndFor
  \State $g_{\max}\gets \max g$;\quad $g_{\min}\gets \min g$
  \State $\mathrm{RCL}\gets\{(h,d)\in\mathcal{C}: g(h,d)\ge g_{\max}-\alpha(g_{\max}-g_{\min})\}$
  \State sorteie $(h^\star,d^\star)$ uniformemente da $\mathrm{RCL}$
  \State $X_{h^\star,d^\star}\mathrel{+}=1$;\quad $v \gets v + g(h^\star,d^\star)$
\EndFor
\State \Return $S$
\end{algorithmic}
\end{algorithm}

\subsection{Fase de busca local}

A solução construída raramente é ótima localmente. A busca local substitui
iterativamente a solução corrente por um vizinho melhor, parando no primeiro
ótimo local. Como o orçamento é fixo, a vizinhança natural é o \textbf{movimento de
um POI}: retirar uma unidade de um par $(h,d)$ ocupado e inseri-la em outro par
$(h',d')$. Usa-se a estratégia \emph{first-improving} (aceita a primeira melhoria
encontrada), recomendada por Resende e Ribeiro por ser mais rápida e menos sujeita
à convergência prematura que a \emph{best-improving}.

\begin{algorithm}[ht]
\caption{Busca local \emph{first-improving} (movimento de 1 POI)}
\label{alg:ls}
\begin{algorithmic}[1]
\State $v \gets f(S)$
\Repeat
  \State \textit{melhorou} $\gets$ \textbf{falso}
  \ForAll{$(h,d)$ com $X_{h,d}>0$ (ordem aleatória)}
    \ForAll{$(h',d') \neq (h,d)$ (ordem aleatória)}
      \State $S' \gets S$ movendo 1 POI de $(h,d)$ para $(h',d')$
      \If{$f(S') > v + \epsilon$}
        \State $S \gets S'$;\ $v \gets f(S')$;\ \textit{melhorou} $\gets$ \textbf{verdadeiro};\ \textbf{break}
      \EndIf
    \EndFor
    \If{\textit{melhorou}} \textbf{break} \EndIf
  \EndFor
\Until{$\neg$\,\textit{melhorou} \textbf{ou} orçamento de avaliações esgotado}
\State \Return $S$
\end{algorithmic}
\end{algorithm}

% ============================================================================
\section{Implementação}
\label{sec:impl}
% ============================================================================

A implementação reside em \texttt{metaheuristics/methods/grasp.py} e expõe a
função \texttt{run\_grasp(context)}, despachada automaticamente pelo otimizador
quando o usuário escolhe a opção \textbf{B}. Esta seção descreve cada parte.

\subsection{Onde o método se encaixa}

O orquestrador \texttt{walk\_meta\_opt} prepara um objeto \texttt{context}
(pacote de dados) e o entrega ao \emph{runner} do método. Os campos relevantes
para o GRASP são:
\begin{itemize}\setlength\itemsep{1pt}
\item \texttt{context.objective\_state\_nd} --- estado numérico pré-compilado
(matriz baseline, índices fonte/alvo, valores $\alpha$, mapeamentos);
\item \texttt{context.budget} --- orçamento $B$ de POIs;
\item \texttt{context.seeds} --- lista de sementes (define o multi-start);
\item \texttt{context.source\_hex\_ids} e \texttt{context.dimensions} --- formam o
conjunto base $E$.
\end{itemize}

\subsection{Representação numérica e propagação}

Toda a avaliação opera em \texttt{ndarray}, sem \texttt{pandas} no laço quente.
A solução é a \texttt{candidate\_matrix} $X$ de forma $(n,m)$; o estado
pré-compilado \texttt{objective\_state\_nd} guarda \texttt{baseline\_matrix}
($M^0$), os arrays \texttt{source\_indices}, \texttt{target\_indices},
\texttt{alpha\_values} (que representam o conjunto $P$ e os $\alpha_{s,t}$) e o
mapa \texttt{candidate\_to\_indicator\_indices} (que liga cada coluna de $X$ à
coluna correspondente em $M$). A propagação~\eqref{eq:propagacao} é uma
soma dispersa (\emph{scatter-add}) vetorizada.

\subsection{Função de avaliação \texttt{\_evaluate}}

Encapsula a Eq.~\eqref{eq:propagacao} seguida da
Eq.~\eqref{eq:objetivo}: dada $X$, propaga o impacto e devolve $\sum_i\mathrm{IQC}_i$.

\begin{lstlisting}[style=py,caption={Avaliação de um candidato: propaga o impacto e retorna a soma do IQC.}]
def _evaluate(candidate_matrix, objective_state):
    """Score a candidate matrix: propagate impacts then return sum(IQC)."""
    final_indicator_matrix = build_final_indicator_matrix_nd(
        candidate_matrix=candidate_matrix,
        objective_state=objective_state,
    )
    return objective_function(final_indicator_matrix=final_indicator_matrix)['objective_value']
\end{lstlisting}

\subsection{Fase de construção \texttt{\_greedy\_randomized\_construction}}

Implementa o Algoritmo~\ref{alg:cons}. Pontos a destacar:
\begin{itemize}\setlength\itemsep{1pt}
\item começa de $X=\mathbf{0}$ e insere 1 POI por passo, até $B$;
\item por passo, amostra \texttt{sample\_size} pares (\emph{sampled greedy}) e mede
o ganho de cada um chamando \texttt{\_evaluate} com o POI somado e depois revertido;
\item monta a RCL pela Eq.~\eqref{eq:rcl} e sorteia um índice dela;
\item a atualização \texttt{current\_value += gains[chosen]} é \emph{exata}, pois
\texttt{gains[chosen]} já é a diferença entre o valor avaliado e o valor corrente.
\end{itemize}

\begin{lstlisting}[style=py,caption={Construção gulosa-aleatória com RCL amostrada.}]
def _greedy_randomized_construction(objective_state, elements, budget,
                                    alpha, sample_size, rng):
    n_hex = len(objective_state.h3_ids)
    n_dims = len(objective_state.candidate_dimensions)
    candidate = np.zeros((n_hex, n_dims), dtype=np.float64)
    current_value = _evaluate(candidate, objective_state)

    for _ in range(budget):
        if sample_size < len(elements):
            sampled = rng.sample(elements, sample_size)
        else:
            sampled = list(elements)

        gains = np.empty(len(sampled), dtype=np.float64)
        for i, (row, col) in enumerate(sampled):
            candidate[row, col] += 1.0
            gains[i] = _evaluate(candidate, objective_state) - current_value
            candidate[row, col] -= 1.0

        g_max = float(gains.max())
        g_min = float(gains.min())
        threshold = g_max - alpha * (g_max - g_min)        # RCL de maximizacao
        rcl = [i for i in range(len(sampled)) if gains[i] >= threshold]

        chosen = rng.choice(rcl)
        row, col = sampled[chosen]
        candidate[row, col] += 1.0
        current_value += float(gains[chosen])

    return candidate, current_value
\end{lstlisting}

\subsection{Fase de busca local \texttt{\_local\_search}}

Implementa o Algoritmo~\ref{alg:ls} com \emph{first-improving}. A vizinhança é o
movimento de 1 POI; a varredura é aleatorizada (ordem dos pares ocupados e dos
destinos) e limitada por um orçamento de avaliações \texttt{max\_evals}, que
garante tempo de execução previsível mesmo quando a função objetivo é cara.

\begin{lstlisting}[style=py,caption={Busca local first-improving por movimento de POI.}]
def _local_search(candidate, objective_state, elements, rng, max_evals):
    current_value = _evaluate(candidate, objective_state)
    evals = 0
    improved = True

    while improved and evals < max_evals:
        improved = False
        occupied = [(int(r), int(c)) for r, c in zip(*np.nonzero(candidate))]
        rng.shuffle(occupied)

        for (row, col) in occupied:
            destinations = list(elements)
            rng.shuffle(destinations)

            for (row2, col2) in destinations:
                if row2 == row and col2 == col:
                    continue
                if evals >= max_evals:
                    break
                candidate[row, col] -= 1.0
                candidate[row2, col2] += 1.0
                value = _evaluate(candidate, objective_state)
                evals += 1
                if value > current_value + IMPROVEMENT_EPS:
                    current_value = value
                    improved = True
                    break                       # first-improving: aceita e reinicia
                candidate[row, col] += 1.0       # reverte
                candidate[row2, col2] -= 1.0

            if improved or evals >= max_evals:
                break

    return candidate, current_value
\end{lstlisting}

\subsection{Multi-start e retorno \texttt{run\_grasp}}

O laço principal (Algoritmo~\ref{alg:grasp}) percorre as sementes, executa
construção + busca local por semente e guarda a melhor solução. Antes do laço,
calcula o objetivo \emph{baseline} (matriz $X=\mathbf 0$) para medir o ganho.
O conjunto base é restrito aos hexágonos-fonte que de fato propagam impacto
(\texttt{np.unique(source\_indices)}), evitando avaliar pares inócuos.

\begin{lstlisting}[style=py,caption={Laço multi-start e montagem do resultado (trecho central).}]
elements = [(int(row), col) for row in source_row_indices for col in range(n_dims)]
seeds = context.seeds if context.seeds else [0]
baseline_value = _evaluate(np.zeros((len(objective_state.h3_ids), n_dims)), objective_state)

best_value, best_candidate, best_seed = float('-inf'), None, None
for seed in seeds:
    rng = random.Random(int(seed))
    candidate, _ = _greedy_randomized_construction(
        objective_state, elements, context.budget, RCL_ALPHA,
        CONSTRUCTION_SAMPLE_SIZE, rng)
    candidate, value = _local_search(
        candidate, objective_state, elements, rng, LOCAL_SEARCH_MAX_EVALS)
    if value > best_value:
        best_value, best_candidate, best_seed = value, candidate.copy(), int(seed)
\end{lstlisting}

O dicionário de retorno traz \texttt{best\_objective\_value} e um
\texttt{best\_solution\_summary} com a semente vencedora, o \texttt{baseline\_objective},
o \texttt{improvement} e a \texttt{allocation} final --- a lista de itens
$\{$\texttt{h3\_id}, \texttt{dimension}, \texttt{quantity}$\}$ obtida convertendo a
matriz $X$ de volta para forma legível (\texttt{\_candidate\_matrix\_to\_items}).

\subsection{Paralelismo e cache}
\label{sec:paralelismo}

O laço multi-start é \emph{embaraçosamente paralelo}: cada semente executa uma
construção e uma busca local independentes, sem estado compartilhado. Por isso
despachamos \textbf{uma semente por processo}, distribuídas em
$\max(\text{núcleos}-1,\,1)$ \emph{workers} (via \texttt{ProcessPoolExecutor}).
Essa é a estratégia de paralelização independente do GRASP, que Resende e Ribeiro
reportam levar a \emph{speedups} próximos do linear.

Três cuidados garantem corretude e robustez:
\begin{itemize}\setlength\itemsep{1pt}
\item \textbf{Determinismo.} Cada semente possui o seu próprio gerador
\texttt{random.Random(seed)}, de modo que o resultado independe do número de
\emph{workers} e da ordem de conclusão. Os resultados são reordenados por semente
antes de montar as estatísticas, então os arquivos de saída são idênticos em
execução serial ou paralela (verificado por testes de igualdade
serial~$=$~paralelo).
\item \textbf{Sem \emph{oversubscription}.} Como cada \emph{worker} já é um
processo, fixamos o BLAS do \texttt{numpy} em uma \emph{thread} por \emph{worker}
(variáveis \texttt{OMP\_NUM\_THREADS} etc.), evitando contenção de CPU.
\item \textbf{Tolerância a falhas.} Se o ambiente não permitir a criação de
processos, o método captura a exceção e executa \emph{serialmente}, sem
interromper a otimização.
\end{itemize}

Adicionalmente, cada \emph{worker} mantém um \textbf{cache} em memória (por semente)
que memoiza a avaliação pela \emph{conteúdo} da matriz candidata
(\texttt{candidate.tobytes()}), evitando recomputar matrizes idênticas. Como a
função objetivo é não-separável (CRITIC global), matrizes candidatas raramente se
repetem dentro da construção; o cache, portanto, ajuda sobretudo na busca local e
\emph{nunca} prejudica (apenas memoiza), tendo um teto de tamanho para não esgotar
a memória. O ganho de desempenho principal vem do paralelismo.

\subsection{Parâmetros e decisões de projeto}

As escolhas adotadas e seus controles (constantes no topo do módulo) estão na
Tabela~\ref{tab:params}. As três primeiras refletem as decisões feitas: $\alpha$
fixo, busca local \emph{first-improving} e \emph{sampled greedy}. Os dois freios de
custo (\texttt{CONSTRUCTION\_SAMPLE\_SIZE} e \texttt{LOCAL\_SEARCH\_MAX\_EVALS})
existem porque a função objetivo recalcula o CRITIC globalmente a cada avaliação;
as quatro últimas controlam o paralelismo e o cache (Seção~\ref{sec:paralelismo}).

\begin{table}[ht]
\centering
\caption{Parâmetros da implementação.}
\label{tab:params}
\begin{tabular}{@{}lll@{}}
\toprule
\textbf{Constante} & \textbf{Valor} & \textbf{Papel} \\
\midrule
\texttt{RCL\_ALPHA}              & $0{,}2$ & limiar da RCL (Eq.~\ref{eq:rcl}) \\
\texttt{CONSTRUCTION\_SAMPLE\_SIZE} & $64$ & pares avaliados por passo (\emph{sampled greedy}) \\
\texttt{LOCAL\_SEARCH\_MAX\_EVALS}  & $4000$ & teto de avaliações da busca local por semente \\
\texttt{IMPROVEMENT\_EPS}        & $10^{-9}$ & ganho mínimo para aceitar melhoria \\
\texttt{GRASP\_PARALLEL}         & \texttt{True} & liga o multi-start paralelo \\
\texttt{GRASP\_MAX\_WORKERS}     & \texttt{None} & \texttt{None}~$=$~núcleos$-1$; ou inteiro fixo \\
\texttt{GRASP\_EVAL\_CACHE}      & \texttt{True} & memoização da avaliação por \emph{worker} \\
\texttt{GRASP\_EVAL\_CACHE\_MAXSIZE} & $5\times10^{5}$ & teto de entradas do cache por semente \\
\bottomrule
\end{tabular}
\end{table}

% ============================================================================
\section{Execução ilustrada passo a passo}
\label{sec:exemplo}
% ============================================================================

Para tornar a mecânica concreta, executamos o GRASP numa \textbf{instância
didática} reduzida (números reais produzidos pelo próprio código). Para facilitar
a visualização usamos $3$ indicadores e $2$ dimensões alocáveis; a implementação
real usa $11$ indicadores e $7$ dimensões, mas a lógica é idêntica.

\subsection{Instância}

\textbf{Hexágonos:} $\hex{A},\hex{B},\hex{C},\hex{D}$. \quad
\textbf{Indicadores:} \texttt{S\_saude}, \texttt{T\_transporte}, \texttt{I\_seguranca}.\\
\textbf{Dimensões alocáveis:} \texttt{S\_saude} (coluna 0) e
\texttt{T\_transporte} (coluna 1). \quad \textbf{Orçamento:} $B=3$.\\
\textbf{Propagação:} cada fonte $s$ impacta a si mesma ($\alpha=1{,}0$) e o
vizinho seguinte ($\alpha=0{,}5$), no ciclo
$\hex{A}\!\to\!\hex{B}\!\to\!\hex{C}\!\to\!\hex{D}\!\to\!\hex{A}$.

\begin{table}[ht]
\centering
\caption{Matriz \emph{baseline} de indicadores $M^0$. Objetivo baseline
$f(\mathbf 0)=1{,}8389$.}
\label{tab:baseline}
\begin{tabular}{@{}c ccc@{}}
\toprule
Hex & \texttt{S\_saude} & \texttt{T\_transporte} & \texttt{I\_seguranca} \\
\midrule
\hex{A} & 2 & 5 & 3 \\
\hex{B} & 4 & 1 & 6 \\
\hex{C} & 1 & 3 & 2 \\
\hex{D} & 5 & 2 & 4 \\
\bottomrule
\end{tabular}
\end{table}

\subsection{Construção --- passo a passo}

A cada passo mostramos o ganho de \emph{todos} os 8 elementos (aqui
\texttt{sample\_size} cobre o conjunto base inteiro, então não há amostragem),
a RCL (\colorbox{cellrcl}{azul}) e o elemento sorteado
(\colorbox{cellsel}{laranja}). À direita, a matriz de alocação $X$ resultante,
com a célula alterada em \colorbox{cellchg}{verde}.

% -------------------- PASSO 1 --------------------
\paragraph{Passo 1.}
$g_{\max}=0{,}1400$,\ $g_{\min}=-0{,}1572$,\ limiar
$=0{,}1400-0{,}2\cdot(0{,}1400+0{,}1572)=0{,}0806$.

\begin{minipage}[t]{0.55\textwidth}
\centering
\small
\begin{tabular}{@{}lc@{}}
\toprule
Elemento $(h,d)$ & ganho $g$ \\
\midrule
\rowcolor{cellsel}(\hex{A},\,\texttt{S\_saude})      & $+0{,}1400$ \\
(\hex{A},\,\texttt{T\_transporte})                   & $-0{,}1297$ \\
(\hex{B},\,\texttt{S\_saude})                        & $+0{,}0378$ \\
(\hex{B},\,\texttt{T\_transporte})                   & $-0{,}1061$ \\
(\hex{C},\,\texttt{S\_saude})                        & $-0{,}1572$ \\
\rowcolor{cellrcl}(\hex{C},\,\texttt{T\_transporte}) & $+0{,}1287$ \\
(\hex{D},\,\texttt{S\_saude})                        & $-0{,}0325$ \\
(\hex{D},\,\texttt{T\_transporte})                   & $+0{,}0569$ \\
\bottomrule
\end{tabular}
\end{minipage}
\hfill
\begin{minipage}[t]{0.40\textwidth}
\centering
\small
\vspace{2pt}
$X$ após o passo 1:\\[4pt]
\begin{tabular}{@{}c cc@{}}
\toprule
 & \texttt{S\_saude} & \texttt{T\_transp.} \\
\midrule
\hex{A} & \cellcolor{cellchg}1 & 0 \\
\hex{B} & 0 & 0 \\
\hex{C} & 0 & 0 \\
\hex{D} & 0 & 0 \\
\bottomrule
\end{tabular}\\[4pt]
$f = 1{,}9789$
\end{minipage}

\medskip
A RCL contém os dois elementos acima do limiar; o sorteio escolheu
(\hex{A},\,\texttt{S\_saude}).

% -------------------- PASSO 2 --------------------
\paragraph{Passo 2.}
$g_{\max}=0{,}1714$,\ $g_{\min}=-0{,}1531$,\ limiar $=0{,}1065$.
Note como os ganhos \emph{mudaram} em relação ao passo 1 --- é o caráter
adaptativo: a solução parcial alterou o CRITIC e a normalização.

\begin{minipage}[t]{0.55\textwidth}
\centering
\small
\begin{tabular}{@{}lc@{}}
\toprule
Elemento $(h,d)$ & ganho $g$ \\
\midrule
\rowcolor{cellsel}(\hex{A},\,\texttt{S\_saude})      & $+0{,}1714$ \\
(\hex{A},\,\texttt{T\_transporte})                   & $-0{,}1531$ \\
(\hex{B},\,\texttt{S\_saude})                        & $-0{,}0723$ \\
(\hex{B},\,\texttt{T\_transporte})                   & $-0{,}1355$ \\
(\hex{C},\,\texttt{S\_saude})                        & $-0{,}1385$ \\
\rowcolor{cellrcl}(\hex{C},\,\texttt{T\_transporte}) & $+0{,}1626$ \\
(\hex{D},\,\texttt{S\_saude})                        & $-0{,}0583$ \\
(\hex{D},\,\texttt{T\_transporte})                   & $+0{,}0661$ \\
\bottomrule
\end{tabular}
\end{minipage}
\hfill
\begin{minipage}[t]{0.40\textwidth}
\centering
\small
\vspace{2pt}
$X$ após o passo 2:\\[4pt]
\begin{tabular}{@{}c cc@{}}
\toprule
 & \texttt{S\_saude} & \texttt{T\_transp.} \\
\midrule
\hex{A} & \cellcolor{cellchg}2 & 0 \\
\hex{B} & 0 & 0 \\
\hex{C} & 0 & 0 \\
\hex{D} & 0 & 0 \\
\bottomrule
\end{tabular}\\[4pt]
$f = 2{,}1503$
\end{minipage}

\medskip
Sorteio escolheu novamente (\hex{A},\,\texttt{S\_saude}): agora $X_{\hex{A},\,\texttt{S\_saude}}=2$.

% -------------------- PASSO 3 --------------------
\paragraph{Passo 3.}
$g_{\max}=0{,}1786$,\ $g_{\min}=-0{,}2037$,\ limiar $=0{,}1021$.
O ganho de mais um POI de saúde em \hex{A} caiu para $+0{,}0269$ (rendimentos
decrescentes), e agora só (\hex{C},\,\texttt{T\_transporte}) supera o limiar.

\begin{minipage}[t]{0.55\textwidth}
\centering
\small
\begin{tabular}{@{}lc@{}}
\toprule
Elemento $(h,d)$ & ganho $g$ \\
\midrule
(\hex{A},\,\texttt{S\_saude})                        & $+0{,}0269$ \\
(\hex{A},\,\texttt{T\_transporte})                   & $-0{,}1654$ \\
(\hex{B},\,\texttt{S\_saude})                        & $-0{,}2037$ \\
(\hex{B},\,\texttt{T\_transporte})                   & $-0{,}1495$ \\
(\hex{C},\,\texttt{S\_saude})                        & $-0{,}1446$ \\
\rowcolor{cellsel}(\hex{C},\,\texttt{T\_transporte}) & $+0{,}1786$ \\
(\hex{D},\,\texttt{S\_saude})                        & $-0{,}1017$ \\
(\hex{D},\,\texttt{T\_transporte})                   & $+0{,}0658$ \\
\bottomrule
\end{tabular}
\end{minipage}
\hfill
\begin{minipage}[t]{0.40\textwidth}
\centering
\small
\vspace{2pt}
$X$ após o passo 3 (final):\\[4pt]
\begin{tabular}{@{}c cc@{}}
\toprule
 & \texttt{S\_saude} & \texttt{T\_transp.} \\
\midrule
\hex{A} & 2 & 0 \\
\hex{B} & 0 & 0 \\
\hex{C} & 0 & \cellcolor{cellchg}1 \\
\hex{D} & 0 & 0 \\
\bottomrule
\end{tabular}\\[4pt]
$f = 2{,}3289$
\end{minipage}

\subsection{Propagação do impacto}

A solução construída aloca $2$ POIs de \texttt{S\_saude} na fonte \hex{A} e
$1$ POI de \texttt{T\_transporte} na fonte \hex{C}. Aplicando a
Eq.~\eqref{eq:propagacao}:

\begin{itemize}\setlength\itemsep{1pt}
\item \textbf{\hex{A} (2 POIs de saúde):} $\hex{A}\!\to\!\hex{A}$ com
$\alpha=1{,}0$ soma $+2{,}0$ à saúde de \hex{A};
$\hex{A}\!\to\!\hex{B}$ com $\alpha=0{,}5$ soma $+1{,}0$ à saúde de \hex{B}.
\item \textbf{\hex{C} (1 POI de transporte):} $\hex{C}\!\to\!\hex{C}$ com
$\alpha=1{,}0$ soma $+1{,}0$ ao transporte de \hex{C};
$\hex{C}\!\to\!\hex{D}$ com $\alpha=0{,}5$ soma $+0{,}5$ ao transporte de \hex{D}.
\end{itemize}

\begin{table}[ht]
\centering
\caption{Matriz final $M = M^0 + \Delta$. Em \colorbox{cellchg}{verde}, as
células alteradas pela propagação.}
\label{tab:final}
\begin{tabular}{@{}c ccc@{}}
\toprule
Hex & \texttt{S\_saude} & \texttt{T\_transporte} & \texttt{I\_seguranca} \\
\midrule
\hex{A} & \cellcolor{cellchg}4 \;\scriptsize($2{+}2$) & 5 & 3 \\
\hex{B} & \cellcolor{cellchg}5 \;\scriptsize($4{+}1$) & 1 & 6 \\
\hex{C} & 1 & \cellcolor{cellchg}4 \;\scriptsize($3{+}1$) & 2 \\
\hex{D} & 5 & \cellcolor{cellchg}2{,}5 \;\scriptsize($2{+}0{,}5$) & 4 \\
\bottomrule
\end{tabular}
\end{table}

\subsection{Recálculo do CRITIC e do IQC}

Sobre a matriz final, a função objetivo recalcula os pesos
CRITIC~\eqref{eq:critic} e o IQC~\eqref{eq:iqc}:
\begin{equation*}
w = \bigl(\underbrace{0{,}4436}_{\texttt{S\_saude}},\
          \underbrace{0{,}3608}_{\texttt{T\_transporte}},\
          \underbrace{0{,}1956}_{\texttt{I\_seguranca}}\bigr).
\end{equation*}

\begin{table}[ht]
\centering
\caption{IQC por hexágono após a construção e objetivo global.}
\label{tab:iqc}
\begin{tabular}{@{}c c@{}}
\toprule
Hex & $\mathrm{IQC}_i$ \\
\midrule
\hex{A} & 0{,}7424 \\
\hex{B} & 0{,}6392 \\
\hex{C} & 0{,}2706 \\
\hex{D} & 0{,}6767 \\
\midrule
$f(X)=\sum_i \mathrm{IQC}_i$ & \textbf{2{,}3289} \\
\bottomrule
\end{tabular}
\end{table}

\subsection{Resultado e leitura}

Partindo do baseline $f(\mathbf 0)=1{,}8389$, a fase de construção já elevou o
objetivo para $f(X)=2{,}3289$, um \textbf{ganho de $+0{,}4900$}. A busca local
parte daqui tentando movimentos de 1 POI; nesta instância mínima a solução
construída já é ótima local. Em instâncias reais, é a combinação
\emph{construção (diversificação via RCL) $+$ busca local (intensificação) $+$
multi-start} que produz soluções de alta qualidade, exatamente como descrito por
Resende e Ribeiro.

A leitura urbanística do resultado é direta: concentrar POIs de saúde no hexágono
\hex{A} (que também irriga \hex{B}) e um POI de transporte em \hex{C} (que
alcança \hex{D}) foi a alocação de melhor retorno de IQC sob o orçamento dado.

% ============================================================================
\section{Arquivos gerados para análise estatística}
\label{sec:arquivos}
% ============================================================================

Toda execução do otimizador grava automaticamente um conjunto de arquivos
projetados para análise estatística reprodutível e publicação. Esta seção
descreve cada arquivo, o que ele representa, como pode ser usado, quais métodos
estatísticos são adequados e o significado de \emph{cada} atributo (coluna).

\subsection{Organização: duas camadas}

Os arquivos ficam em \texttt{data/optimization\_runs/} e se organizam em duas
camadas complementares (Figura~\ref{fig:tree}):

\begin{enumerate}
\item \textbf{Arquivos-mestre acumulativos} (na raiz): cada execução
\emph{acrescenta} linhas. São os arquivos lidos para análises \emph{entre}
execuções (comparar métodos, perfis, localidades).
\item \textbf{Pasta por execução}
(\texttt{<key\_location>/<perfil>/<método>/<timestamp>\_\_<id>/}): guarda o detalhe
completo e reprodutível de uma única execução.
\end{enumerate}

\begin{figure}[ht]
\centering
\begin{lstlisting}[basicstyle=\ttfamily\footnotesize,frame=single,
  backgroundcolor=\color{codebg},rulecolor=\color{gray!40}]
data/optimization_runs/
|-- all_runs_summary.csv        (mestre: 1 linha por execucao)
|-- all_per_seed.csv            (mestre: 1 linha por execucao x semente)
|-- all_per_hexagon_iqc.csv     (mestre: 1 linha por execucao x hexagono)
|-- all_allocations.csv         (mestre: 1 linha por execucao x POI alocado)
`-- <key_location>/<perfil>/<metodo>/<timestamp>__<id>/
    |-- run_metadata.json          metadados completos (reprodutibilidade)
    |-- run_report.txt             resumo legivel
    |-- run_summary.csv            resumo de 1 linha (esta execucao)
    |-- per_seed.csv               resultado por semente
    |-- per_hexagon_iqc.csv        IQC baseline vs otimizado por hexagono
    |-- per_hexagon_indicators.csv indicadores (formato longo)
    |-- allocation.csv             melhor alocacao
    |-- critic_weights.csv         pesos CRITIC baseline vs otimizado
    `-- convergence.csv            objetivo por passo de construcao
\end{lstlisting}
\caption{Estrutura de diretórios dos arquivos de saída.}
\label{fig:tree}
\end{figure}

\subsection{Identificadores comuns (chaves de junção)}

As quatro tabelas-mestre compartilham um conjunto de colunas de identificação,
que servem de \emph{chaves} para cruzamentos (\texttt{merge}/\texttt{join}) e
agrupamentos (\texttt{group\_by}). Elas aparecem em todos os arquivos longos e não
são repetidas nas tabelas de atributos seguintes.

\begin{center}\small
\begin{tabular}{@{}>{\ttfamily}l l p{8.4cm}@{}}
\toprule
\normalfont Atributo & \normalfont Tipo & \normalfont Descrição \\
\midrule
run\_id          & texto    & Identificador único da execução (UUID curto). Chave primária que liga \emph{todos} os arquivos de uma mesma execução. \\
location         & texto    & Nome da localidade (legível). \\
key\_location    & texto    & Chave saneada da localidade (usada em caminhos). \\
walking\_profile & texto    & Perfil de caminhada (\texttt{average\_adult}, \texttt{elderly}, \texttt{athlete}). \\
method\_code     & texto    & Código do método (\texttt{B} = GRASP). \\
\bottomrule
\end{tabular}
\end{center}

% ----------------------------------------------------------------------------
\subsection{\texttt{all\_runs\_summary.csv} \,/\, \texttt{run\_summary.csv}}
% ----------------------------------------------------------------------------

\textbf{O que representa.} Uma linha por execução, com todos os indicadores
agregados. É o arquivo \emph{central} para estatística: cada linha é uma
observação experimental (uma combinação localidade $\times$ perfil $\times$ método
$\times$ conjunto de sementes).

\textbf{Como usar e quais métodos aplicar.}
\begin{itemize}\setlength\itemsep{1pt}
\item \emph{Comparar métodos/perfis/localidades:} agrupar por fator e comparar
\texttt{best\_objective} ou \texttt{improvement\_pct}. Testar normalidade
(Shapiro--Wilk) e homocedasticidade (Levene); se atendidas, \textbf{ANOVA
fatorial} com pós-teste de \textbf{Tukey HSD}; caso contrário, \textbf{Kruskal--Wallis}
(ou \textbf{Friedman}, se pareado pela mesma instância localidade$\times$perfil)
com pós-teste de \textbf{Dunn/Nemenyi}. Reportar tamanho de efeito
($\eta^2$, $\epsilon^2$ ou $\delta$ de Cliff).
\item \emph{Modelar o ganho:} regressão linear (ou modelo misto)
$\texttt{improvement\_pct} \sim \text{método}+\text{perfil}+\text{localidade}+\texttt{budget}$.
\item \emph{Robustez do método:} usar \texttt{obj\_seed\_std} / \texttt{obj\_seed\_mean}
(coeficiente de variação) como medida de estabilidade entre sementes.
\end{itemize}

\paragraph{Atributos --- identificação e estado.}
\begin{center}\small
\begin{tabular}{@{}>{\ttfamily}l l p{8.0cm}@{}}
\toprule
\normalfont Atributo & \normalfont Tipo & \normalfont Descrição \\
\midrule
run\_id        & texto      & Identificador único da execução. \\
timestamp\_utc & ISO-8601   & Data/hora de início em UTC. \\
location, key\_location, walking\_profile & texto & Ver chaves comuns. \\
method\_code, method\_name & texto & Código e nome do método. \\
status         & texto      & \texttt{ok}, \texttt{error}, \texttt{placeholder}, etc. \\
run\_dir       & texto      & Caminho da pasta detalhada desta execução. \\
\bottomrule
\end{tabular}
\end{center}

\paragraph{Atributos --- configuração.}
\begin{center}\small
\begin{tabular}{@{}>{\ttfamily}l l p{8.0cm}@{}}
\toprule
\normalfont Atributo & \normalfont Tipo & \normalfont Descrição \\
\midrule
budget                 & inteiro & Número $B$ de POIs alocados. \\
n\_seeds               & inteiro & Número de sementes (reinícios multi-start). \\
best\_seed             & inteiro & Semente que produziu a melhor solução. \\
objective\_metric      & texto   & Nome da métrica (\texttt{sum\_iqc}). \\
optimization\_direction& texto   & Direção (\texttt{maximize}). \\
rcl\_alpha             & real    & Parâmetro $\alpha$ da RCL. \\
construction\_sample\_size & inteiro & Pares avaliados por passo (\emph{sampled greedy}). \\
local\_search\_max\_evals & inteiro & Teto de avaliações da busca local. \\
h3\_resolution         & inteiro & Resolução H3 da malha. \\
distance               & inteiro & Distância (m) usada na geração dos dados. \\
\bottomrule
\end{tabular}
\end{center}

\paragraph{Atributos --- objetivo e ganho.}
\begin{center}\small
\begin{tabular}{@{}>{\ttfamily}l l p{8.0cm}@{}}
\toprule
\normalfont Atributo & \normalfont Tipo & \normalfont Descrição \\
\midrule
baseline\_objective & real & $f(\mathbf 0)$: soma do IQC sem intervenção. \\
best\_objective     & real & $f(X^\star)$: soma do IQC da melhor solução. \\
improvement         & real & $f(X^\star)-f(\mathbf 0)$ (ganho absoluto). \\
improvement\_pct    & real & Ganho relativo ($100\cdot$ \texttt{improvement}/\texttt{baseline}). \\
\bottomrule
\end{tabular}
\end{center}

\paragraph{Atributos --- distribuição entre sementes.}
Resumo do objetivo \emph{após busca local} sobre as \texttt{n\_seeds} execuções
(útil para caracterizar a variabilidade estocástica do GRASP).
\begin{center}\small
\begin{tabular}{@{}>{\ttfamily}l l p{8.0cm}@{}}
\toprule
\normalfont Atributo & \normalfont Tipo & \normalfont Descrição \\
\midrule
obj\_seed\_mean & real & Média do objetivo entre as sementes. \\
obj\_seed\_std  & real & Desvio-padrão amostral (estabilidade do método). \\
obj\_seed\_min  & real & Pior objetivo entre as sementes. \\
obj\_seed\_max  & real & Melhor objetivo ($=$ \texttt{best\_objective}). \\
\bottomrule
\end{tabular}
\end{center}

\paragraph{Atributos --- distribuição do IQC por hexágono.}
Estatísticas descritivas do IQC \emph{entre hexágonos} (não entre sementes).
Respondem diretamente à pergunta ``como o IQC varia nesta localidade/perfil''.
\begin{center}\small
\begin{tabular}{@{}>{\ttfamily}l l p{7.8cm}@{}}
\toprule
\normalfont Atributo & \normalfont Tipo & \normalfont Descrição \\
\midrule
iqc\_base\_mean, iqc\_base\_std & real & Média e desvio do IQC \emph{baseline} entre hexágonos. \\
iqc\_opt\_mean, iqc\_opt\_std   & real & Média e desvio do IQC \emph{otimizado}. \\
iqc\_opt\_min, iqc\_opt\_median, iqc\_opt\_max & real & Mínimo, mediana e máximo do IQC otimizado. \\
delta\_iqc\_mean, delta\_iqc\_std & real & Média e desvio da variação $\Delta\mathrm{IQC}$ por hexágono. \\
delta\_iqc\_min, delta\_iqc\_max & real & Maior queda e maior aumento de IQC entre hexágonos. \\
delta\_iqc\_positive\_count & inteiro & Quantos hexágonos tiveram o IQC \emph{aumentado}. \\
\bottomrule
\end{tabular}
\end{center}

\paragraph{Atributos --- instância, solução e tempo.}
\begin{center}\small
\begin{tabular}{@{}>{\ttfamily}l l p{8.0cm}@{}}
\toprule
\normalfont Atributo & \normalfont Tipo & \normalfont Descrição \\
\midrule
n\_hexagons        & inteiro & Total de hexágonos da malha. \\
n\_source\_hexagons& inteiro & Hexágonos-fonte elegíveis (que propagam impacto). \\
n\_candidate\_dimensions & inteiro & Dimensões de POI alocáveis. \\
allocated\_pois    & inteiro & POIs alocados na melhor solução ($=B$). \\
distinct\_pairs    & inteiro & Pares (hex, dimensão) distintos usados. \\
runtime\_seconds\_total & real & Tempo total da execução (todas as sementes). \\
\bottomrule
\end{tabular}
\end{center}

% ----------------------------------------------------------------------------
\subsection{\texttt{all\_per\_seed.csv} \,/\, \texttt{per\_seed.csv}}
% ----------------------------------------------------------------------------

\textbf{O que representa.} Uma linha por semente. Expõe a \emph{variabilidade
estocástica} interna: como cada reinício se saiu na construção e após a busca
local.

\textbf{Como usar e quais métodos aplicar.}
\begin{itemize}\setlength\itemsep{1pt}
\item \emph{Variabilidade:} média, desvio, coeficiente de variação e intervalos de
confiança \textbf{bootstrap} de \texttt{local\_search\_objective}.
\item \emph{Contribuição da busca local:} testar se \texttt{local\_search\_gain}
difere de zero (teste $t$ pareado ou \textbf{Wilcoxon} de postos sinalizados,
comparando construção vs.\ pós-busca-local).
\item \emph{Eficiência:} relacionar \texttt{runtime\_seconds} com a qualidade
(análise tempo$\times$qualidade), base para \emph{time-to-target}.
\end{itemize}

\begin{center}\small
\begin{tabular}{@{}>{\ttfamily}l l p{8.0cm}@{}}
\toprule
\normalfont Atributo & \normalfont Tipo & \normalfont Descrição \\
\midrule
seed                    & inteiro & Semente do gerador pseudoaleatório. \\
construction\_objective & real & Objetivo ao fim da fase de construção. \\
local\_search\_objective& real & Objetivo após a busca local (qualidade final desta semente). \\
local\_search\_gain     & real & Ganho da busca local (\texttt{local\_search} $-$ \texttt{construction}). \\
improvement\_over\_baseline & real & \texttt{local\_search\_objective} $-$ \texttt{baseline}. \\
runtime\_seconds        & real & Tempo desta semente (construção $+$ busca local). \\
distinct\_pairs         & inteiro & Pares (hex, dimensão) distintos na solução desta semente. \\
\bottomrule
\end{tabular}
\end{center}

% ----------------------------------------------------------------------------
\subsection{\texttt{all\_per\_hexagon\_iqc.csv} \,/\, \texttt{per\_hexagon\_iqc.csv}}
% ----------------------------------------------------------------------------

\textbf{O que representa.} Uma linha por hexágono (georreferenciada), com o IQC
antes e depois da intervenção. É a base para a \emph{análise espacial} dos
ganhos.

\textbf{Como usar e quais métodos aplicar.}
\begin{itemize}\setlength\itemsep{1pt}
\item \emph{Distribuição:} histogramas e \emph{boxplots} de \texttt{delta\_iqc}
por localidade/perfil; teste de Wilcoxon pareado entre \texttt{iqc\_baseline} e
\texttt{iqc\_optimized}.
\item \emph{Autocorrelação espacial:} \textbf{I de Moran} (global) e \textbf{LISA}
(local) sobre \texttt{delta\_iqc} usando \texttt{latitude}/\texttt{longitude} para
a matriz de vizinhança --- verifica se os ganhos formam aglomerados espaciais.
\item \emph{Pontos quentes:} \textbf{Getis--Ord $G_i^\star$} para identificar
\emph{hot/cold spots} de melhoria.
\item \emph{Validação:} comparar \texttt{iqc\_baseline} (recalculado pela função
objetivo) com \texttt{iqc\_pipeline} (valor original do pipeline) como checagem de
consistência.
\end{itemize}

\begin{center}\small
\begin{tabular}{@{}>{\ttfamily}l l p{8.0cm}@{}}
\toprule
\normalfont Atributo & \normalfont Tipo & \normalfont Descrição \\
\midrule
h3\_id        & texto & Identificador H3 do hexágono. \\
latitude, longitude & real & Coordenadas do centro do hexágono. \\
iqc\_baseline & real & IQC do hexágono \emph{sem} intervenção (métrica da função objetivo). \\
iqc\_optimized& real & IQC do hexágono na melhor solução. \\
delta\_iqc    & real & Variação $\texttt{iqc\_optimized}-\texttt{iqc\_baseline}$. \\
iqc\_pipeline & real & IQC original calculado pelo pipeline (referência). \\
\bottomrule
\end{tabular}
\end{center}

% ----------------------------------------------------------------------------
\subsection{\texttt{per\_hexagon\_indicators.csv}}
% ----------------------------------------------------------------------------

\textbf{O que representa.} Formato \emph{longo} (\emph{tidy}): uma linha por
hexágono $\times$ indicador, com o valor antes/depois. Permite ver \emph{quais
indicadores} (saúde, transporte, etc.) moveram o IQC e \emph{onde}.

\textbf{Como usar e quais métodos aplicar.} Decompor o ganho por dimensão;
\emph{heatmaps} indicador$\times$hexágono; correlação entre \texttt{delta} dos
indicadores e \texttt{delta\_iqc}; análise de contribuição (qual indicador explica
melhor a variação do IQC).

\begin{center}\small
\begin{tabular}{@{}>{\ttfamily}l l p{8.0cm}@{}}
\toprule
\normalfont Atributo & \normalfont Tipo & \normalfont Descrição \\
\midrule
h3\_id          & texto & Hexágono. \\
indicator       & texto & Nome do indicador (ex.\ \texttt{S\_saude}). \\
value\_baseline & real  & Valor do indicador antes da intervenção. \\
value\_optimized& real  & Valor após a propagação dos POIs alocados. \\
delta           & real  & \texttt{value\_optimized} $-$ \texttt{value\_baseline}. \\
\bottomrule
\end{tabular}
\end{center}

% ----------------------------------------------------------------------------
\subsection{\texttt{all\_allocations.csv} \,/\, \texttt{allocation.csv}}
% ----------------------------------------------------------------------------

\textbf{O que representa.} Uma linha por item alocado: \emph{onde} e \emph{quanto}
de cada dimensão foi inserido na melhor solução.

\textbf{Como usar e quais métodos aplicar.} Frequência de alocação por dimensão e
por hexágono; mapas de calor da alocação; relação entre alocação e
\texttt{delta\_iqc} (cruzando por \texttt{h3\_id}); padrões de alocação entre
perfis (ex.\ idosos vs.\ atletas priorizam dimensões diferentes?).

\begin{center}\small
\begin{tabular}{@{}>{\ttfamily}l l p{8.0cm}@{}}
\toprule
\normalfont Atributo & \normalfont Tipo & \normalfont Descrição \\
\midrule
h3\_id    & texto   & Hexágono-fonte que recebeu os POIs. \\
dimension & texto   & Dimensão do POI alocado. \\
quantity  & inteiro & Quantidade de POIs alocados naquele par. \\
latitude, longitude & real & Coordenadas do hexágono-fonte. \\
\bottomrule
\end{tabular}
\end{center}

% ----------------------------------------------------------------------------
\subsection{\texttt{critic\_weights.csv}}
% ----------------------------------------------------------------------------

\textbf{O que representa.} Os pesos CRITIC de cada indicador, \emph{antes} e
\emph{depois} da intervenção. Como o CRITIC é orientado a dados, a alocação altera
a importância relativa dos indicadores.

\textbf{Como usar e quais métodos aplicar.} Comparar \texttt{weight\_baseline} vs.\
\texttt{weight\_optimized} (gráfico de barras pareado); análise de sensibilidade do
IQC aos pesos; discutir como a intervenção reequilibra a importância das
dimensões.

\begin{center}\small
\begin{tabular}{@{}>{\ttfamily}l l p{8.0cm}@{}}
\toprule
\normalfont Atributo & \normalfont Tipo & \normalfont Descrição \\
\midrule
indicator        & texto & Nome do indicador. \\
weight\_baseline & real  & Peso CRITIC no cenário baseline (soma $=1$). \\
weight\_optimized& real  & Peso CRITIC na melhor solução (soma $=1$). \\
\bottomrule
\end{tabular}
\end{center}

% ----------------------------------------------------------------------------
\subsection{\texttt{convergence.csv}}
% ----------------------------------------------------------------------------

\textbf{O que representa.} O objetivo registrado \emph{a cada passo} da fase de
construção, para \emph{cada} semente (formato longo). Suporta curvas de
convergência e gráficos \emph{time-to-target} (TTT), padrão na literatura do GRASP.

\textbf{Como usar e quais métodos aplicar.} Curvas objetivo$\times$passo com
banda de confiança (média $\pm$ IC entre sementes); \emph{área sob a curva} como
medida de rapidez de convergência; \textbf{TTT plots} (probabilidade acumulada de
atingir um alvo) com ajuste de distribuição \textbf{exponencial deslocada}
(Aiex, Resende e Ribeiro).

\begin{center}\small
\begin{tabular}{@{}>{\ttfamily}l l p{8.0cm}@{}}
\toprule
\normalfont Atributo & \normalfont Tipo & \normalfont Descrição \\
\midrule
seed      & inteiro & Semente correspondente à trajetória. \\
step      & inteiro & Passo de construção ($1,\dots,B$). \\
objective & real    & Valor do objetivo após inserir o $k$-ésimo POI. \\
\bottomrule
\end{tabular}
\end{center}

% ----------------------------------------------------------------------------
\subsection{\texttt{run\_metadata.json}}
% ----------------------------------------------------------------------------

\textbf{O que representa.} Os metadados completos da execução, em JSON, garantindo
\textbf{reprodutibilidade} (requisito para periódicos A1/A2). Reúne, num só
arquivo: identificação, configuração, objetivo, \texttt{parameters} (todos os
parâmetros do método), \texttt{instance} (tamanho da instância), \texttt{seeds}
(lista completa) e \texttt{environment} --- versões de Python, \texttt{numpy},
\texttt{pandas}, plataforma e o \emph{commit} do Git. Com ele, qualquer execução
pode ser refeita exatamente.

\textbf{Como usar.} É o registro canônico de \emph{proveniência}: antes de
reprocessar ou comparar resultados, confira \texttt{environment.git\_commit} (versão
exata do código), \texttt{seeds} (sementes que reproduzem o multi-start) e
\texttt{input\_files} (dados de origem). Diferentemente dos CSVs, é hierárquico:
campos escalares no topo e três blocos aninhados (\texttt{parameters},
\texttt{instance}, \texttt{environment}), detalhados a seguir.

\paragraph{Campos de topo.}
\begin{center}\small
\begin{tabular}{@{}>{\ttfamily}l l p{7.8cm}@{}}
\toprule
\normalfont Campo & \normalfont Tipo & \normalfont Descrição \\
\midrule
run\_id        & texto    & Identificador único da execução (liga ao \texttt{run\_id} dos CSVs). \\
timestamp\_utc & ISO-8601 & Data/hora de início em UTC (\texttt{datetime.isoformat}). \\
location, key\_location, walking\_profile & texto & Ver chaves comuns. \\
method\_code, method\_name & texto & Código e nome do método. \\
status         & texto    & Estado final (\texttt{ok}, \texttt{error}, \dots). \\
budget         & inteiro  & Orçamento $B$ de POIs. \\
seeds          & lista[int] & \emph{Todas} as sementes usadas (não só a melhor) --- reproduz o multi-start. \\
n\_seeds       & inteiro  & Tamanho de \texttt{seeds}. \\
objective\_metric & texto & Métrica otimizada (\texttt{sum\_iqc}). \\
optimization\_direction & texto & Direção (\texttt{maximize}). \\
baseline\_objective & real & $f(\mathbf 0)$: objetivo sem intervenção. \\
best\_objective & real    & $f(X^\star)$: objetivo da melhor solução. \\
improvement, improvement\_pct & real & Ganho absoluto e relativo (\%) sobre o baseline. \\
best\_seed     & inteiro  & Semente que produziu a melhor solução. \\
runtime\_seconds\_total & real & Tempo total (todas as sementes), em segundos. \\
h3\_resolution & inteiro  & Resolução H3 da malha. \\
distance       & inteiro  & Distância (m) usada na geração dos dados. \\
message        & texto    & Resumo textual do resultado retornado pelo método. \\
input\_files   & objeto   & Mapa \texttt{nome}$\to$\texttt{caminho} dos arquivos de entrada (rastreabilidade da origem dos dados). \\
parameters     & objeto   & Parâmetros do método (Tabela abaixo). \\
instance       & objeto   & Tamanho da instância (Tabela abaixo). \\
environment    & objeto   & Ambiente de execução (Tabela abaixo). \\
\bottomrule
\end{tabular}
\end{center}

\paragraph{Bloco \texttt{parameters}.} Os controles efetivamente usados na
execução (constantes do topo de \texttt{grasp.py}; ver Tabela~\ref{tab:params}).
\begin{center}\small
\begin{tabular}{@{}>{\ttfamily}l l p{7.8cm}@{}}
\toprule
\normalfont Campo & \normalfont Tipo & \normalfont Descrição \\
\midrule
rcl\_alpha                 & real      & Limiar $\alpha$ da RCL (Eq.~\ref{eq:rcl}). \\
construction\_sample\_size & inteiro   & Pares avaliados por passo (\emph{sampled greedy}). \\
local\_search\_max\_evals  & inteiro   & Teto de avaliações da busca local por semente. \\
improvement\_eps           & real      & Ganho mínimo ($10^{-9}$) para aceitar uma melhoria na busca local. \\
parallel\_workers          & inteiro   & Número de processos \emph{worker} efetivamente usados (1 $=$ serial). \\
eval\_cache                & booleano  & Se a memoização da avaliação por \emph{worker} esteve ativa. \\
\bottomrule
\end{tabular}
\end{center}

\paragraph{Bloco \texttt{instance}.} Descreve o tamanho do problema resolvido.
\begin{center}\small
\begin{tabular}{@{}>{\ttfamily}l l p{7.8cm}@{}}
\toprule
\normalfont Campo & \normalfont Tipo & \normalfont Descrição \\
\midrule
n\_hexagons          & inteiro     & Total de hexágonos da malha ($n$). \\
n\_source\_hexagons  & inteiro     & Hexágonos-fonte elegíveis a receber POIs (que propagam impacto). \\
n\_candidate\_dimensions & inteiro & Número $m$ de dimensões de POI alocáveis. \\
candidate\_dimensions & lista[texto] & Nomes das dimensões alocáveis (define as colunas de $X$). \\
\bottomrule
\end{tabular}
\end{center}

\paragraph{Bloco \texttt{environment}.} Captura o ambiente para reprodução exata.
\begin{center}\small
\begin{tabular}{@{}>{\ttfamily}l l p{7.8cm}@{}}
\toprule
\normalfont Campo & \normalfont Tipo & \normalfont Descrição \\
\midrule
python      & texto & Versão do interpretador Python. \\
numpy       & texto & Versão do \texttt{numpy}. \\
pandas      & texto & Versão do \texttt{pandas}. \\
platform    & texto & Descrição da plataforma/sistema operacional. \\
git\_commit & texto & \emph{Hash} curto do \emph{commit} do código (ou \texttt{null} fora de um repositório Git). \\
\bottomrule
\end{tabular}
\end{center}

% ----------------------------------------------------------------------------
\subsection{\texttt{run\_report.txt}}
% ----------------------------------------------------------------------------

\textbf{O que representa.} Um resumo legível por humanos da execução (mesma
informação do \texttt{run\_summary.csv}, formatada). Serve para conferência rápida
e para anexos do relatório/artigo.

% ----------------------------------------------------------------------------
\subsection{Receituário de análises}
% ----------------------------------------------------------------------------

A Tabela~\ref{tab:receituario} mapeia perguntas de pesquisa típicas aos arquivos e
métodos correspondentes.

\begin{table}[ht]
\centering
\caption{Da pergunta de pesquisa ao arquivo e ao método estatístico.}
\label{tab:receituario}
\small
\begin{tabular}{@{}p{4.6cm} p{4.2cm} p{5.6cm}@{}}
\toprule
\textbf{Pergunta} & \textbf{Arquivo(s)} & \textbf{Método} \\
\midrule
Um método supera os outros? &
\texttt{all\_runs\_summary} &
Friedman + pós-teste de Nemenyi (pareado por instância); ou ANOVA + Tukey. \\
\addlinespace
O ganho varia por perfil/localidade? &
\texttt{all\_runs\_summary} &
ANOVA fatorial / modelo misto; $\eta^2$ como efeito. \\
\addlinespace
Quão estável é o método? &
\texttt{all\_per\_seed} &
Coeficiente de variação; IC bootstrap; Wilcoxon (construção vs.\ busca local). \\
\addlinespace
Os ganhos de IQC formam aglomerados no espaço? &
\texttt{all\_per\_hexagon\_iqc} &
I de Moran (global), LISA (local), Getis--Ord $G_i^\star$. \\
\addlinespace
Quais indicadores explicam o ganho? &
\texttt{per\_hexagon\_indicators} &
Correlação \texttt{delta}$\leftrightarrow$\texttt{delta\_iqc}; heatmaps; contribuição. \\
\addlinespace
Onde os POIs são alocados? &
\texttt{all\_allocations} &
Frequências; mapas de calor; comparação entre perfis. \\
\addlinespace
A intervenção reequilibra os pesos? &
\texttt{critic\_weights} &
Comparação pareada baseline vs.\ otimizado; sensibilidade. \\
\addlinespace
Quão rápido o método converge? &
\texttt{convergence} &
Curvas média $\pm$ IC; TTT plots (exponencial deslocada). \\
\bottomrule
\end{tabular}
\end{table}

\paragraph{Ponto de partida prático.} Carregar \texttt{all\_runs\_summary.csv} e
cruzar com os demais mestres pela coluna \texttt{run\_id} cobre a maior parte das
análises. Os arquivos por execução entram quando se deseja o detalhe fino (curvas
de convergência, mapas espaciais, decomposição por indicador).

% ============================================================================
\section{Considerações finais e extensões}
% ============================================================================

A implementação cobre o GRASP \emph{tradicional} de duas fases. A estrutura
modular (avaliação isolada em \texttt{\_evaluate}; construção, busca local e
multi-start independentes) facilita evoluções descritas por Resende e Ribeiro:

\begin{itemize}\setlength\itemsep{2pt}
\item \textbf{Reactive GRASP}: auto-ajuste do parâmetro $\alpha$ ao longo das
iterações, em vez de fixo, aumentando a robustez;
\item \textbf{Path-Relinking}: intensificação que explora trajetórias entre a
solução corrente e um conjunto de soluções de elite (corresponde ao método
híbrido da opção E);
\item \textbf{Estratégias de reinício} (\emph{restart}) para reduzir a cauda da
distribuição de tempo até alvo.
\end{itemize}

O principal desafio prático do problema --- a não-separabilidade da função
objetivo (CRITIC global) --- foi tratado com \emph{sampled greedy} na construção e
um orçamento de avaliações na busca local, mantendo o tempo de execução
previsível sem comprometer a qualidade.

% ============================================================================
\begin{thebibliography}{9}
\bibitem{feo1995}
T.~A.~Feo and M.~G.~C.~Resende.
\emph{Greedy randomized adaptive search procedures}.
Journal of Global Optimization, 6:109--133, 1995.

\bibitem{rr2016}
M.~G.~C.~Resende and C.~C.~Ribeiro.
\emph{Optimization by GRASP: Greedy Randomized Adaptive Search Procedures}.
Springer, 2016.

\bibitem{rr2019}
M.~G.~C.~Resende and C.~C.~Ribeiro.
\emph{Greedy Randomized Adaptive Search Procedures: Advances and Extensions}.
In: Handbook of Metaheuristics, cap.~6, Springer, 2019.
\end{thebibliography}

\end{document}
